\section{Conclusion}
\label{sec:conclusion}

MRI-ReportGenerator is an end-to-end, application-positioned pipeline that turns a single sagittal T2
cervical MRI into a structured, literature-grounded, no-diagnosis report. Its contribution is disciplined
engineering rather than a new algorithm: composing three validated segmentation engines, organizing
measurements into five clinically meaningful groups, grounding every threshold in a cited primary source
with explicit modality caveats, and --- most importantly --- validating the assembled system honestly.

Because no public dataset pairs cervical MRI with per-case expert measurements, we adopted a
threshold-crossing design built on the principle that measurements are disease-agnostic geometric/signal
detectors anchored on healthy norms. On real cervical MRIs, the canal/cord group separated strongly
($p=0.0001$) and the vertebral compression screen was correctly silent on a non-compression cohort. Two
leads were then \emph{revised} by larger, design-appropriate cohorts: the SPINEPS Cobb method reads
correct lordosis but, on a representative multi-site healthy sample, does not discriminate myelopathy (a
validated measurement, not a screen); and the disc group, after the harness detected and we fixed a
sign-level bug, separates only modestly through disc geometric spread (AUC $0.62$), with disc signal and
bulge documented as non-discriminators. Four measurement defects were found and fixed and the
assessement layer was wired end-to-end --- the value being not a clean scoreboard but a validation
disciplined enough to overturn its own encouraging early numbers.

The honest limitations are clear: small, partly selected samples; a single (severe) disease spectrum; no
per-case ground truth; an under-tested compression-fracture arm; a disc group only partially validated;
and an assessement layer whose per-case clinical accuracy still awaits an institutional read. The path
forward is correspondingly concrete: scale Group~3 to a random MMCSD draw; pursue an institutional read
(e.g.\ AUBMC) for true per-case accuracy and for the demographic-adjusted thresholds; secure a labelled
cervical compression-fracture set for the one untested axis; and complete the deployed demonstration.

The broader lesson is methodological. In a domain without ground truth, the most valuable engineering is
not a clever model but a transparent chain of evidence --- cited thresholds, documented mistakes, kept
negative results, precision distinguished from accuracy, and a validation that doubles as a bug detector.
That chain, more than any single number, is what makes an automated clinical measurement tool trustworthy
enough to put in front of a physician for review.
